\section{模型评价与结论}

\subsection{结果是否满足题目要求}

本项目使用 UCF101 官方 split 1，在 3783 个视频的官方测试集上完成了三条路线、
三种采样和等成本时间空间配置。核心配置均使用三个训练种子，给出均值与标准差。
主实验的 TSN 线性头、TSN 时序注意力头和 R(2+1)D-18 分别相对 R3D-18 提升
4.93、6.34 和 7.60 个百分点，均达到题目要求的至少 4 个百分点。

显存方面，单个训练步骤和正式特征提取均执行 2048 MiB 硬门禁。代表性技术消融的
最高 reserved 峰值为 950 MiB，正式时间优先时序头的峰值为 1912 MiB，未超过
2048 MiB。因此，从数据、精度目标、多种子统计和训练预算四个核心维度看，本研究
完成了题目的可验收要求。

\subsection{主要结论}

第一，低显存条件下轻量时空分解卷积具有很高价值。R(2+1)D-18 达到最高绝对精度
83.97\%，同时特征提取 reserved 峰值只有 218 MiB。它的代价是吞吐低于二维路线，
适合优先考虑精度且能够接受较高离线计算成本的场景。

第二，二维时序路线在效率与精度之间更均衡。TSN-2D 的特征提取吞吐为 52.98
video/s，高于 R(2+1)D-18 的 22.20 video/s，Top-1 达到 82.70\%。对于流式或
接近实时的应用，二维路线更合适。

第三，均匀分段优于连续稠密。相同时序头下，连续稠密只有 78.60\%，比均匀采样
低 4.10 个百分点。运动自适应达到 82.02\%，但额外解码和评分使吞吐下降约 68\%，
没有带来正收益。

第四，在等帧像素成本下应优先空间分辨率。空间优先达到 84.19\%，时间优先只有
61.71\%，差距为 22.48 个百分点。时间维并非不重要，时序注意力相对平均池化
仍有约 1.41 个百分点收益，但增加帧数必须满足基本的空间细节条件。

\subsection{模型优点}

\begin{enumerate}
  \item 使用官方切分和 group 隔离，测试集只用于最终评估。
  \item 数据、代码、训练种子、显存峰值和结果文件均可追溯。
  \item 路线、采样和输入预算分离控制，避免把架构变化与采样变化混杂。
  \item 主体实验在真实 4 GiB 显卡上完成，2 GiB 门禁不是离线估算。
  \item 同时报告精度、标准差、显存、吞吐和自适应预处理成本。
\end{enumerate}

\subsection{局限性与改进方向}

主实验冻结 Kinetics-400 预训练主干，只训练分类头或时序聚合头，因此结论适用于
“冻结主干迁移训练”，不等同于全网络端到端微调。若拥有更大显存，可进一步比较
最后阶段的局部微调和全网络微调。

当前主输入为 8 帧，动作阶段划分仍是离散采样。自适应方法只使用灰度帧差，未来可
改用轻量光流、帧级显著性或多尺度候选窗口。三维路线虽然精度最高，但当前批次
较小，吞吐仍受 R(2+1)D-18 的计算量限制。后续可尝试时序通道剪枝和知识蒸馏。

错误归因显示人体动作组和细粒度周期动作仍是主要短板。可引入局部时序卷积、
轨迹描述子和更强的空间细粒度特征。若扩展 RGB 与光流双流，应继续维持相同显存
上限并单列光流预处理成本。

\subsection{推广}

本文的“双峰值显存门禁、固定输入成本、同种子配对比较、额外预处理成本单列”
框架可推广到其他资源受限视频任务，例如课堂行为识别、工业操作合规检测、
体育动作分析、老人跌倒监测和家庭服务机器人。对于边缘设备，二维路线适合在线
推断，轻量三维路线适合事件完成后的高精度复核。
